% Options for packages loaded elsewhere
\PassOptionsToPackage{unicode}{hyperref}
\PassOptionsToPackage{hyphens}{url}
\documentclass[
  ignorenonframetext,
]{beamer}
\newif\ifbibliography
\usepackage{pgfpages}
\setbeamertemplate{caption}[numbered]
\setbeamertemplate{caption label separator}{: }
\setbeamercolor{caption name}{fg=normal text.fg}
\beamertemplatenavigationsymbolsempty
% remove section numbering
\setbeamertemplate{part page}{
  \centering
  \begin{beamercolorbox}[sep=16pt,center]{part title}
    \usebeamerfont{part title}\insertpart\par
  \end{beamercolorbox}
}
\setbeamertemplate{section page}{
  \centering
  \begin{beamercolorbox}[sep=12pt,center]{section title}
    \usebeamerfont{section title}\insertsection\par
  \end{beamercolorbox}
}
\setbeamertemplate{subsection page}{
  \centering
  \begin{beamercolorbox}[sep=8pt,center]{subsection title}
    \usebeamerfont{subsection title}\insertsubsection\par
  \end{beamercolorbox}
}
% Prevent slide breaks in the middle of a paragraph
\widowpenalties 1 10000
\raggedbottom
\AtBeginPart{
  \frame{\partpage}
}
\AtBeginSection{
  \ifbibliography
  \else
    \frame{\sectionpage}
  \fi
}
\AtBeginSubsection{
  \frame{\subsectionpage}
}
\usepackage{iftex}
\ifPDFTeX
  \usepackage[T1]{fontenc}
  \usepackage[utf8]{inputenc}
  \usepackage{textcomp} % provide euro and other symbols
\else % if luatex or xetex
  \usepackage{unicode-math} % this also loads fontspec
  \defaultfontfeatures{Scale=MatchLowercase}
  \defaultfontfeatures[\rmfamily]{Ligatures=TeX,Scale=1}
\fi
\usepackage{lmodern}
\ifPDFTeX\else
  % xetex/luatex font selection
\fi
% Use upquote if available, for straight quotes in verbatim environments
\IfFileExists{upquote.sty}{\usepackage{upquote}}{}
\IfFileExists{microtype.sty}{% use microtype if available
  \usepackage[]{microtype}
  \UseMicrotypeSet[protrusion]{basicmath} % disable protrusion for tt fonts
}{}
\makeatletter
\@ifundefined{KOMAClassName}{% if non-KOMA class
  \IfFileExists{parskip.sty}{%
    \usepackage{parskip}
  }{% else
    \setlength{\parindent}{0pt}
    \setlength{\parskip}{6pt plus 2pt minus 1pt}}
}{% if KOMA class
  \KOMAoptions{parskip=half}}
\makeatother
\setlength{\emergencystretch}{3em} % prevent overfull lines
\providecommand{\tightlist}{%
  \setlength{\itemsep}{0pt}\setlength{\parskip}{0pt}}
% Slide layout defaults for warning-clean scientific decks.
\usepackage{etoolbox}
\IfFileExists{xurl.sty}{\usepackage{xurl}}{}
\IfFileExists{seqsplit.sty}{\usepackage{seqsplit}}{\newcommand{\seqsplit}[1]{#1}}
\protected\def\breakseq#1{\seqsplit{#1}}
\protected\def\breaktt#1{\begingroup\ttfamily\seqsplit{#1}\endgroup}
\setlength{\emergencystretch}{6em}
\tolerance=5000
\hbadness=10000
\hfuzz=1pt
\setlength{\tabcolsep}{2pt}
\AtBeginEnvironment{longtable}{\tiny\renewcommand{\arraystretch}{0.86}\setlength{\tabcolsep}{1pt}}
\AtBeginEnvironment{tabular}{\tiny\renewcommand{\arraystretch}{0.86}\setlength{\tabcolsep}{1pt}}

% Natbib and cross-reference fallbacks — slides don't load natbib
% or cleveref, but combined-PDF manuscript prose may emit these
% commands. The fallback renders citations as a bracketed key list
% and cross-references as detokenized labels so slides don't fail on
% undefined control sequences or raw underscores. \providecommand is
% a no-op if packages load later.
\providecommand{\citep}[1]{[#1]}
\providecommand{\citet}[1]{#1}
\providecommand{\citealp}[1]{#1}
\providecommand{\citeauthor}[1]{#1}
\providecommand{\citeyear}[1]{#1}
\providecommand{\cref}[1]{\texttt{\detokenize{#1}}}
\providecommand{\Cref}[1]{\texttt{\detokenize{#1}}}

\newtheorem{proposition}{Proposition}
\newtheorem{hypothesis}{Hypothesis}
\usepackage{bookmark}
\IfFileExists{xurl.sty}{\usepackage{xurl}}{} % add URL line breaks if available
\urlstyle{same}
% fallback for those not using the hyperref driver hyperxmp:
\makeatletter
\@ifundefined{xmpquote}{\newcommand{\xmpquote}[1]{#1}}{}
\makeatother
\hypersetup{
  hidelinks,
  pdfcreator={LaTeX via pandoc}}

\author{\texorpdfstring{}{}}
\date{}

\begin{document}

\section{Discussion: Limits, Failure Modes, and Future
Replacement}\label{sec:discussion}

\begin{frame}[allowframebreaks]{Discussion: Limits, Failure Modes, and
Future Replacement}
DigiPPPiP belongs to the broader landscape of relational technology, but
it takes a specific position within it. Many contemporary systems
replace, simulate, or optimize social contact. DigiPPPiP instead uses
technology to give two human partners another surface for mutual
attention. That difference matters ethically: the goal is not
frictionless engagement, retention, or automated companionship, but a
bounded practice in which each partner can see and respond to the
other's agency.

The strongest scholarship added by the digital turn is not a single
theory but a triangulation. CSCW explains why the shared drawing space
is an interactional resource, with social translucence and workspace
awareness providing sharper design constraints for what partner activity
should become visible \breakseq{[@tang1991collaborativework;
@erickson2000socialtranslucence; @gutwin2002workspaceawareness]}.
Distance and grounding research explain why remote drawing needs
evidence of understanding, repair, and re-entry rather than only a
persistent file \breakseq{[@olson2000distance; @clark1991grounding]}. Presence
and embodiment research explain why lightweight traces can carry felt
co-presence without recreating the whole body \breakseq{[@lombard1997presence;
@biocca1997cyborg; @lee2004presence]}. Digital-intimacy and
relatedness-technology research explains why rituals, minimal signals,
traces, and joint action matter for strong-tie relationships while also
warning that many prototypes lack robust effect evidence
\breakseq{[@kaye2006clicked; @hassenzahl2012love; @neustaedter2012intimacy;
@vetere2005mediating; @mcveighschultz2015couple;
@wenhart2025relatedness]}.

The same triangulation clarifies the risks. Persistent canvases can
support shared memory only when contextual privacy, collective
information practice, values-in-design, and dyadic boundary rules are
explicit \breakseq{[@nissenbaum2011contextualprivacy; @palen2003privacy;
@dourish2006collectiveprivacy; @petronio2020cpm; @shilton2012values]}.
Emerging work on AI-mediated collaborative drawing in romantic
relationships is relevant only if AI remains a reflective mediator
rather than the relationship partner or the author of the shared gesture
\breakseq{[@won2026venus]}. Companion-AI scholarship makes that boundary
stronger: systems designed to satisfy companionship needs can raise
replacement, social motivation, and social-skill concerns for human
relationships, even when some users experience benefits
\breakseq{[@malfacini2025companionai]}. Mixed-initiative and human-AI
interaction scholarship therefore set a high bar: optionality, control,
reversibility, disclosure, and recovery from wrong suggestions are not
interface polish but conditions for preserving partner agency
\breakseq{[@horvitz1999mixedinitiative; @deterding2017mixedinitiative;
@amershi2019humanai]}. Hyperscanning and active inference supply
testable measurement and modeling languages, but not automatic causal
proof \breakseq{[@friston2017process; @hamilton2021hyperscanning]}.

The most important revision to the project brief is evidentiary
discipline. The framework draws from hyperscanning, active inference,
digital art therapy, neuroergonomics, phenomenology, relational
aesthetics, and digital placemaking, but those literatures do not yet
prove DigiPPPiP's outcomes. The manuscript therefore phrases claims as
design rationales and empirical hypotheses unless a cited source
directly supports the narrower statement. The source-quality map in
\breakseq{[@fig:source\_quality\_map]} makes this discipline visible in the
artifact itself.

The computational model also changes the manuscript's posture. Rather
than leaving abstract phrases such as ``dyadic generative model'' or
``narrative entropy'' as metaphors, the project implements small
deterministic primitives and keeps them in
\breakseq{[@sec:formalisms\_appendix]}. Those primitives are intentionally
modest. Their value is that they expose assumptions, make figures
reproducible, and prevent quantitative claims from being invented in
prose.

The first-principles review leaves only a small set of hard constraints:
two people, a shared field of marks, perceivable agency, temporal
structure, and accountable persistence. Most other elements are soft
constraints or hypotheses. Digital capture, neural measurement, AI
assistance, place tagging, and clinical translation are optional
branches that must justify themselves against that small core. This
makes the framework easier to test and harder to inflate.

The systems review adds a second discipline: every optional branch needs
a feedback loop, a causal assumption, an ethics gate, and a reversal
path. Without those pieces, the branch is not mature enough to carry
manuscript claims. This is especially important for long-distance use,
because a persistent surface can become useful, irrelevant, burdensome,
or intrusive depending on the dyad's actual uptake. It is also important
for AI assistance, where reversibility and authorship clarity matter
more than novelty.

The same restraint applies to health, intimacy, and accessibility.
DigiPPPiP may become useful in dyadic digital-health or creative-arts
contexts, but the current evidence supports only a research agenda and
design rationale. Claims about treatment, relationship improvement,
universal access, AI-supported insight, privacy-safe archives, or
community-scale placemaking require study designs that directly measure
those outcomes. The manuscript should therefore remain hospitable to
future empirical replacement: the conceptual figures can be swapped for
real data, but the claim-strength rules should stay.

The textosexual insight from PPPiP remains useful: relationships can be
strengthened through shared marks, shared texts, and shared visual
artifacts \breakseq{[@mikhailova2018pppip]}. DigiPPPiP updates that insight for
networked surfaces. The mark can now be simultaneous or delayed, local
or remote, physical or digital, private or archived. The relational
requirement is unchanged: partners must attend, respond, and co-author.
If future studies show that ordinary shared drawing, video chat, or
another low-friction activity produces the same outcomes with less
complexity, DigiPPPiP should narrow to the contexts where its
persistence, accessibility, place responsiveness, or event logging add
specific value.
\end{frame}

\end{document}
